%% ===================================================================
%% One-Venue DEX Settlement on Lux
%% Companion paper to LP-311 (lps/LP-311-one-venue-dex-settlement.md)
%% and LP-032 (DEX VM). Standalone; compiles with shared/luxcover.
%% ===================================================================

\documentclass[11pt,a4paper]{article}
\usepackage{shared/luxcover}
\usepackage[utf8]{inputenc}
\usepackage[T1]{fontenc}
\usepackage{amsmath,amssymb,amsfonts,amsthm}
\usepackage{graphicx}
\usepackage{hyperref}
\usepackage{booktabs}
\usepackage{listings}
\usepackage{xcolor}
\usepackage{float}
\usepackage{array}
\usepackage{tabularx}
\input{shared/lstlang}

\hypersetup{colorlinks=true,linkcolor=black,citecolor=blue,urlcolor=blue}
\lstset{basicstyle=\ttfamily\scriptsize,breaklines=true,frame=single,numbers=left,numberstyle=\tiny\color{gray}}

\newtheorem{theorem}{Theorem}[section]
\newtheorem{proposition}{Proposition}[section]
\newtheorem{definition}{Definition}[section]
\newtheorem{lemma}{Lemma}[section]
\newtheorem{remark}{Remark}[section]

\title{One-Venue DEX Settlement\\Stateless Conduit, Synchronous Atomic Carry-Fills, and Post-Quantum Cross-Network Transport}
\author{Zach Kelling \\
\textit{Lux Architecture Team}}
\date{June 2026}

\begin{document}
\luxcoverpage

\begin{abstract}
The Lux decentralized exchange is \emph{one venue}---the D-Chain, a virtual machine that runs a
central-limit order book as a consensus-native operation (LP-032)---reached from every EVM chain
through \emph{one settlement primitive}, the \texttt{0x9999} precompile (LP-9999). This paper
specifies and proves correct the settlement \emph{layer} that composes them. We make the
\texttt{0x9999} precompile a \emph{stateless conduit} that holds no order book and no ledger:
the single book and the single ledger live on the venue. We show that settlement on the Lux
primary network is synchronous within a single block---preserving the Uniswap-V4
\texttt{swap()} interface that returns a balance delta in the same transaction---via atomic
shared-memory \emph{carry-fills}, and we prove this construction is consensus-safe whereas the
na\"ive per-validator venue call is not. We show that settlement across networks is asynchronous
via warp/ICM, with the venue's own post-quantum consensus quorum certificate---not a bespoke
exchange certificate---as the sole cross-network trust anchor, and we prove the cross-network
path adds zero trust surface beyond the venue's existing validator set. Finally, we show that
removing state from the conduit dissolves an empty-account reaping hazard that otherwise strands
funds on every EVM chain, and that one venue makes LUX a universal cross-chain fee currency with
shared, unfragmented liquidity.
\end{abstract}

\tableofcontents

\section{Introduction}

A decentralized exchange must answer three questions: \emph{where do orders match}, \emph{where
does custody live}, and \emph{how does value move across chain boundaries}. The Lux design answers
the first two with a single object---the D-Chain venue (LP-032)---and the third with a single
primitive---the \texttt{0x9999} settlement precompile (LP-9999). The discipline that keeps the
system simple is that these answers stay separate: matching and custody are the venue's; movement
is the conduit's; and the only thing that varies between chains is the transport that carries
intents to the venue and fills back.

That discipline had drifted. The \texttt{0x9999} precompile accumulated an in-trie matching engine
and a per-account ledger inside its own account storage, becoming a second venue beside the
D-Chain venue. The drift was not merely inelegant: it manifested as a reproducible,
fund-stranding production failure (Section~\ref{sec:reaping}) that blocked every real fill on every
EVM chain. The failure is the symptom; the disease is the braiding of \emph{where state lives}
into the settlement adapter. This paper decomplects them and proves the resulting system correct.

\paragraph{Contributions.}
\begin{enumerate}
\item A stateless-conduit settlement architecture (Section~\ref{sec:arch}) in which the venue is
      the unique holder of the order book and ledger, the conduit holds only escrow, and transport
      is the single orthogonal axis.
\item A consensus-safety result for synchronous in-block settlement
      (Section~\ref{sec:sync}): the atomic carry-fills construction yields deterministic replay
      (Theorem~\ref{thm:determinism}), while the per-validator venue call admits a safety violation
      (Theorem~\ref{thm:fork}).
\item A zero-new-trust result for cross-network settlement (Section~\ref{sec:async}): a
      cross-network fill is accepted exactly when a $\geq\!\tfrac{2}{3}$-of-venue-stake certificate
      validates it (Theorem~\ref{thm:trust}), reusing the venue's post-quantum consensus
      certificate rather than a bespoke exchange certificate.
\item A reaping-immunity result (Theorem~\ref{thm:reap}) and a universal-fee-rail construction
      (Section~\ref{sec:fee}).
\end{enumerate}

\section{The empty-account reaping hazard}
\label{sec:reaping}

The \texttt{0x9999} precompile is \emph{always-on}: it is dispatchable by its address whether or
not a state account object exists there. When the precompile owns a ledger, it writes that ledger
into its own account's storage. We recall the Ethereum emptiness predicate.

\begin{definition}[EIP-158/161 emptiness]
A state account is \emph{empty} iff $\mathrm{nonce}=0 \wedge \mathrm{balance}=0 \wedge
\mathrm{codeHash}=\texttt{keccak256}(\varepsilon)$. Storage contents are not part of the predicate.
\end{definition}

The per-transaction finalization step \texttt{Finalise}$(\mathit{deleteEmptyObjects}=\texttt{true})$
deletes every empty account \emph{together with all of its storage}. An ERC-20 deposit credits the
token contract's storage (so $\texttt{balanceOf}(\mathit{token},\texttt{0x9999})$ rises) and writes
the precompile's reserve and credit rows into \texttt{0x9999}'s storage---leaving \texttt{0x9999}'s
own account empty (no balance, no nonce, no code). Finalization then reaps it.

A real-\texttt{StateDB} probe matrix (geth host, EIP-158 enabled) isolates the discriminator: the
empty account, not the nested token call.

\begin{table}[H]
\centering
\begin{tabularx}{\textwidth}{Xccc c}
\toprule
\textbf{deposit shape} & \textbf{nested call} & \textbf{\texttt{0x9999} non-empty} &
\textbf{EIP-158} & \textbf{ledger survives \texttt{Finalise(true)}} \\
\midrule
native (adds balance)        & yes & \textbf{yes} & on  & \textbf{yes} \\
ERC-20 (token storage only)  & yes & no          & on  & \textbf{no (reaped)} \\
ERC-20, no nested call       & no  & no          & on  & \textbf{no (reaped)} \\
ERC-20                       & yes & no          & off & yes \\
\bottomrule
\end{tabularx}
\caption{Empty-account reaping isolated. Native deposits survive because the balance credit makes
\texttt{0x9999} non-empty; ERC-20 deposits do not, with or without a nested call; disabling EIP-158
prevents reaping. The discriminator is emptiness, not the nested call.}
\end{table}

\paragraph{Consequence.} No market can be funded, so no fill ever settles
($\mathit{DEXFill}$ count $=0$ on every EVM chain), and any ERC-20 deposited to \texttt{0x9999} is
stranded: its claim row is reaped while the tokens remain in the token contract. Because
\texttt{0x9999} is always-on on every Lux EVM past the activation fork, the exposure is latent on
every chain.

\paragraph{Why the marker is the wrong fix.} Forcing $\mathrm{nonce}(\texttt{0x9999})=1$ keeps the
account non-empty and suppresses the symptom, but it preserves the disease: a precompile owning a
ledger, now defended against the host. The correct remedy is to remove the state from the conduit.
Section~\ref{sec:arch} does so; Theorem~\ref{thm:reap} shows the hazard then cannot arise.

\section{The decomplected architecture}
\label{sec:arch}

\subsection{Values, venue, conduit, transport}

\begin{definition}[Settlement values]
An \emph{order intent} is the tuple a trader submits: market, side, limit, size, and an escrow
reference; it is content-addressed and chain-agnostic. A \emph{fill receipt} is the tuple the venue
produces: the matched legs as per-leg asset deltas
$\langle \mathit{rail}, \mathit{owner}, \mathit{asset}, \mathit{amount}, \mathit{spent}\rangle$
together with a quorum certificate; it is self-authenticating and chain-agnostic.
\end{definition}

\begin{definition}[Venue]
The \emph{venue} is the D-Chain (LP-032): the unique holder of the order book
($\texttt{order:}, \texttt{book:}, \texttt{market:}, \texttt{trade:}$), the ledger
($\texttt{balance:}, \texttt{locked:}$), and the seam-escrow record
($\texttt{seamintent:}, \texttt{seamout:}$). Matching is a VM-native operation at block execution,
deterministic and parallel across markets.
\end{definition}

\begin{definition}[Conduit]
The \emph{conduit} is the \texttt{0x9999} precompile, identical on every EVM chain. It performs
exactly: escrow the physical asset; forward an intent to the venue; apply a fill the venue
produced; release escrow. It holds no order book, reserve, or credit ledger, and writes no durable
row to its own account storage (Section~\ref{sec:conduit}). It does not branch on the transport.
\end{definition}

\begin{definition}[Transport]
A \emph{transport} provides $\textsc{Submit}(\mathit{intent}) \to \mathit{ref}$ (forward C$\to$D,
revert-safe) and $\textsc{Apply}(\mathit{ref}) \to \mathit{receipt}$ (deliver D$\to$C, value-bound,
exactly-once). Two implementations exist: \emph{atomic} (same network as the venue; synchronous)
and \emph{warp} (a separate network; asynchronous). The venue and the conduit never branch on which.
\end{definition}

A new chain---L1, L2, or L3---adds no settlement code; it binds a transport. This is the sense in
which all chains ``work the same.''

\section{Synchronous settlement on the primary network}
\label{sec:sync}

On the Lux primary network the C-Chain, the D-Chain venue, and any co-resident L1/L2/L3 execute
inside a single \texttt{luxd} process with atomic shared memory between chains. A swap therefore
settles synchronously within one block and returns a balance delta from \texttt{swap()}, preserving
the Uniswap-V4 interface.

\subsection{Carry-fills}

\begin{definition}[Carry-fills]\label{def:carry}
The block proposer computes the match \emph{once} against the venue's consensus-accepted state and
records the resulting fill in the block via atomic shared memory. Validators do not re-run the
match; they replay the carried fill from the block bytes and verify consistency with the venue
accepted state they also hold. The C-Chain block and the corresponding venue state transition are
co-committed within the single-proposer block.
\end{definition}

We model a block's settlement as a function of two inputs: the block bytes $b$ and the venue's
accepted state $\sigma_D$ at the block's parent. Write $\mathsf{Settle}(b,\sigma_D)$ for the
post-state a validator computes.

\begin{theorem}[Synchronous determinism]\label{thm:determinism}
Under carry-fills (Definition~\ref{def:carry}) with (C-S1) a single proposer-computed match,
(C-S2) settlement reading only the venue's accepted state---never live, wall-clock, stream-ordered,
or non-deterministically-accelerated state---and (C-S3) no per-validator venue query during
verification, any two honest validators that agree on $b$ and on $\sigma_D$ compute identical
post-states: $\mathsf{Settle}(b,\sigma_D)=\mathsf{Settle}(b,\sigma_D)$ byte-for-byte.
\end{theorem}

\begin{proof}
The settlement-relevant inputs are exactly $b$ and $\sigma_D$. The block bytes $b$ are agreed by
gossip and consensus on the C-Chain block; $\sigma_D$ is agreed by consensus on the venue (both
chains being co-resident and at accepted state). By (C-S2) $\mathsf{Settle}$ reads no other input;
by (C-S1) the carried fill in $b$ is the unique match, not a per-validator recomputation; by (C-S3)
no validator introduces an independent venue read. Hence $\mathsf{Settle}$ is a pure function of
the two agreed inputs, and pure functions of equal inputs return equal outputs. A GPU accelerator
is admitted only when known-answer-tested byte-identical to the reference matcher, so it does not
enlarge the input set. \qed
\end{proof}

\begin{theorem}[The per-validator venue call forks]\label{thm:fork}
Consider the alternative construction in which each validator, during verification of a swap
transaction, independently queries the venue for a match. There is a reachable execution in which
two honest validators return different match results and therefore disagree on the block's
validity---a safety violation.
\end{theorem}

\begin{proof}[Proof (by the existing adversarial witness)]
Let two validators verify the same swap against venue states that differ in the transient ordering
of a concurrently-relayed order---a difference invisible to consensus because the venue query is
not a function of agreed state alone. One validator matches against a book containing the
concurrent order and returns fill $f_1$; the other matches against a book without it and returns
$f_2 \neq f_1$. The two validators then hold different post-states for the same block and vote
incompatibly. This execution is realized by the D-Chain proxy's adversarial test
\texttt{TestRED\_PerValidatorRelay\_SplitsConsensus}, which exhibits a concrete split. The degree
of freedom is the per-validator venue read; Theorem~\ref{thm:determinism} removes it. \qed
\end{proof}

Theorems~\ref{thm:determinism} and~\ref{thm:fork} together are the standard's mandate: synchronous
settlement \emph{must} be carry-fills, never a Verify-time venue call.

\subsection{Developer interface}
\texttt{swap()} keeps its V4 signature and returns a balance delta in the same transaction---no
intent acknowledgement, no ABI change, no second call. Applications on any primary-network chain
integrate the DEX as they would a local pool.

\section{Asynchronous settlement across networks}
\label{sec:async}

A chain on a separate network---its own network identifier, no atomic shared memory with the
venue---settles across the boundary via warp/ICM. Settlement is asynchronous: the conduit escrows
the input and forwards the intent; the fill arrives later carrying the venue's certificate. The V4
call site is unchanged; the realized fill is read from the D$\to$C settlement.

\subsection{The acceptance predicate and zero new trust}

Let $V_D(h)$ be the venue's canonical validator set at P-Chain height $h$, with stake weights, and
let a certificate $\pi$ over a message $M$ at height $h$ \emph{validate} iff the signing weight is
$\geq \tfrac{2}{3}$ of $\mathrm{totalWeight}(V_D(h))$ under the registry-pinned scheme.

\begin{theorem}[Cross-network zero-new-trust]\label{thm:trust}
The conduit accepts a cross-network fill receipt $r$ iff (1) $r$ is the payload of a warp message
$M$ whose certificate $\pi$ validates against $V_D(h)$ with quorum read from configuration (never
from $M$); (2) $M.\mathit{sourceChainID}$ is the venue and $M.\mathit{networkID}$ is local; (3) the
reference is unconsumed and is marked consumed before credit; and (4) the decoded receipt satisfies
the settlement binds. Under these rules, an adversary that controls the transport (the relayer) and
$<\tfrac{1}{3}$ of the venue's stake cannot cause acceptance of any fill the venue did not produce.
The cross-network trust surface equals $V_D$.
\end{theorem}

\begin{proof}
Acceptance requires a validating $\pi$ (1). By the unforgeability of the aggregate scheme and
quorum being read from configuration with a hard $\geq\!\tfrac{2}{3}$ floor, producing a validating
$\pi$ requires signatures of $\geq\tfrac{2}{3}$ of $V_D(h)$'s weight. An adversary with
$<\tfrac{1}{3}$ of venue stake cannot reach that threshold even colluding with all non-venue
parties, since the floor is on venue weight and is not attacker-settable. The relayer's control is
confined to delivery: it cannot alter the payload (the certificate covers the message bytes),
cannot replay (3), and cannot redirect (2). Hence acceptance implies $\geq\tfrac{2}{3}$ of venue
stake attested $r$, i.e.\ the venue produced it. The only trust assumption invoked is honesty of
$\geq\tfrac{2}{3}$ of $V_D$---which is exactly the trust required to trade on the venue at all.
Therefore no trust beyond $V_D$ is added. \qed
\end{proof}

\paragraph{No bespoke exchange certificate.} A DEX-specific aggregate with a verifier registry was
prototyped and deprecated after adversarial review found a value-blind certificate rebind and an
attacker-settable wire quorum collapsing to a single signer. The standard reuses the consensus
certificate and never mints a second trust mechanism; condition (1) above pins quorum to
configuration precisely to exclude the wire-quorum attack.

\subsection{Post-quantum certificate}
The venue certificate is the consensus quorum certificate, post-quantum under LP-020/LP-073. The
primary lane is a Corona (Module-LWE) threshold group signature whose validity structurally implies
the $\geq\tfrac{2}{3}$ floor and which reuses the venue's existing consensus group key (no new
exchange DKG, hence no trusted-dealer reintroduction). The fallback lane is an ML-DSA (FIPS-204)
certificate set summed to $\geq\tfrac{2}{3}$ stake at verification, requiring no DKG. A
\texttt{certType} registry selects classical, Corona, or ML-DSA with no ABI change, profile-gated
so classical and post-quantum compose rather than braid. Both lanes verify deterministically and
inline-safely; no non-determinism enters the certificate check, so Theorem~\ref{thm:determinism}
extends to the cross-network apply.

\section{The stateless conduit}
\label{sec:conduit}

The conduit exposes four operations, each escrow-then-transport, with no durable write to its own
account storage:

\begin{table}[H]
\centering
\begin{tabularx}{\textwidth}{lX}
\toprule
\textbf{operation} & \textbf{effect} \\
\midrule
\texttt{deposit} & escrow in (native balance, or ERC-20 into the token's trie) $\to$ \textsc{Submit}(credit) \\
\texttt{swap}    & escrow input leg $\to$ \textsc{Submit}(order); same-network: receive the carried fill and \emph{return the balance delta}; cross-network: deliver the V4 delta when the fill lands \\
\texttt{fill}    & \textsc{Apply}(ref) $\to$ release escrow to the bound owner (value-bound, exactly-once) \\
\texttt{withdraw}& \textsc{Submit}(withdraw) $\to$ venue debits its ledger $\to$ D$\to$C release $\to$ \texttt{fill} \\
\bottomrule
\end{tabularx}
\end{table}

\begin{theorem}[Reaping immunity]\label{thm:reap}
If the conduit's only state mutations per transaction are (a) its native balance, (b) writes to a
token contract's storage, (c) host transport operations on the atomic backend, and (d) a
transaction-scoped reentrancy guard in transient storage, then
$\texttt{Finalise}(\mathit{deleteEmptyObjects}=\texttt{true})$ over the conduit account neither
loses settlement state nor strands escrow.
\end{theorem}

\begin{proof}
No persistent \texttt{SetState} to the conduit account survives the transaction: (a) is the account
balance, not storage; (b) is another account's storage; (c) is the host atomic backend, a distinct
account; (d) is transient storage, discarded at transaction end and never trie-resident. Hence after
the transaction the conduit account has no own storage rows. If the account is empty it is reaped,
but reaping removes only an account with no balance, nonce, code, or storage---so nothing of value
is lost; the ERC-20 escrow is the token's storage (untouched by reaping the conduit), and the
claim on it is a venue ledger row (not on the conduit at all). The precompile remains dispatchable
by address regardless of account existence. Therefore no settlement state is lost and no escrow is
stranded. \qed
\end{proof}

The nonce marker is thus unnecessary and is removed. Custody follows: the tokens are a fact (the
token's storage), the claim is a venue ledger row, and conservation reduces to
$\sum_{\mathit{asset}} \mathrm{escrow}(\mathit{asset}) \geq \sum_{\mathit{asset}}
\mathrm{venue\text{-}backed\text{-}releasable}(\mathit{asset})$, enforced fail-closed by the venue's
debit and by the value-bound conduit release.

\section{Universal LUX fee rail}
\label{sec:fee}

One venue makes LUX a universal cross-chain fee currency with shared liquidity.

\begin{proposition}[Fee-rail liveness]\label{prop:fee}
If every chain lists initial $\langle \mathit{native}, \mathrm{LUX}\rangle$ liquidity in the venue
at genesis, then every transaction can pay a LUX-denominated fee from the chain's native asset by a
single conduit swap, atomic with the transaction, synchronous on the primary network.
\end{proposition}

\begin{proof}
By hypothesis a $\langle \mathit{native}, \mathrm{LUX}\rangle$ market exists in the venue from block
zero, so the venue can match a $\mathit{native}\!\to\!\mathrm{LUX}$ order of the fee size against
genesis liquidity. The conduit escrows the native fee input and forwards the order; on the primary
network the fill is carried in the same block (Definition~\ref{def:carry}), so the LUX proceeds are
available within the transaction to pay the fee. The swap and the fee debit are committed together
with the transaction. \qed
\end{proof}

The fee currency need not be the chain's native asset. We generalise to \emph{any} asset with a
venue route to LUX (the fee-abstraction standard, LP-312): a transaction declares a fee currency
$X$ and a bound $\overline{x}$ on the amount of $X$ it authorises for fees.

\begin{proposition}[Fee abstraction: liveness and conservation]\label{prop:feeabs}
Let a transaction owe a LUX fee $f$ and declare fee currency $X$ with bound $\overline{x}$. On the
primary network, either (i) the venue's accepted-state liquidity admits an exact-output swap
producing $f$ LUX from some $x\le\overline{x}$ of $X$, in which case the fee is funded in-block, the
recipient receives $f$ in LUX, and value is conserved
\[
   \text{payer } -x\,(X),\qquad
   \text{liquidity } +x\,(X)\ -f\,(\mathrm{LUX}),\qquad
   \text{recipient } +f\,(\mathrm{LUX});
\]
or (ii) no such $x\le\overline{x}$ exists, in which case the transaction is rejected before
execution with no state change. No asset is minted; the recipient is never credited notionally.
\end{proposition}

\begin{proof}
The fee swap is an exact-output venue swap for output $f$, bounded input $\overline{x}$, executed at
fee-buy through the stateless conduit (\S\ref{sec:conduit}) and settled by carry-fills
(Def.~\ref{def:carry}); by Thm~\ref{thm:determinism} it is replay-deterministic. In case (i) the
swap fills against accepted liquidity, moving exactly $x$ of $X$ from payer to liquidity and $f$ of
LUX from liquidity to the fee account, which pays the recipient; summing the three legs gives the
stated conservation, and no leg creates value (it is a swap plus the ordinary gas debit). In case
(ii) admissibility is the deterministic predicate ``the exact-output swap settles within
$\overline{x}$''; its failure is detected at fee-buy (the proposer when building, validators when
replaying, identically), and the transaction is invalid with no state change---the
insufficient-funds semantics. The recipient receives LUX from real liquidity in (i) and nothing in
(ii); there is never a notional credit. \qed
\end{proof}

\begin{remark}
Proposition~\ref{prop:feeabs} is the onboarding unlock: a payer holding zero LUX but holding a
bridged DEX-liquid asset (e.g.\ USDT) transacts immediately, the fee converted in-block through the
one venue. Admissibility is permissionless and self-gating---an asset is a valid fee currency
exactly when its swap settles---so no allowlist and no price oracle are introduced, and by
Thm~\ref{thm:trust}/\ref{thm:determinism} no trust beyond the venue is added. The conversion must be
synchronous (a transaction cannot await an asynchronous cross-network settlement to pay its own
fee), so on a sovereign network the fee asset is converted against a local $\langle X,\mathrm{LUX}
\rangle$ pool rather than the remote venue (LP-312~\S8).
\end{remark}

Because all chains settle to the same book, $\langle \mathit{native}, \mathrm{LUX}\rangle$ (and every
$\langle X, \mathrm{LUX}\rangle$) liquidity is shared rather than fragmented per chain---the property
a per-chain in-trie book cannot provide, and a further reason the second book is removed.

\section{Implementation status}

The venue (LP-032) and the always-on conduit are live on mainnet and testnet; the D-Chain matcher
is bootstrapped and serving. No real fill has settled to date because of the reaping hazard of
Section~\ref{sec:reaping}, which is reproduced on a real \texttt{StateDB} harness. The
stateless-conduit decomplect, the carry-fills synchronous path, the warp cross-network transport,
the post-quantum certificate gating, and the fee rail are specified here and in LP-311; the
consensus post-quantum certificate lanes (Corona, ML-DSA) and the D-Chain proxy's atomic relay
primitive exist. Conformance is reached when a real cross-asset fill settles---synchronously on the
primary network, asynchronously across networks---with the conduit holding zero reapable storage,
the venue holding the sole book and ledger, the certificate as the sole cross-network trust, and the
nonce marker absent.

\section{Related work}

LP-032 specifies the D-Chain venue and the fill receipt. LP-9999 specifies the \texttt{0x9999}
precompile ABI. LP-007 and LP-008 specify CLOB matching and multi-asset trading; LP-040/041/042
specify AMM liquidity sources (not matchers). LP-020 and LP-073 specify the Quasar consensus
certificate and Pulsar post-quantum finality reused for cross-network settlement. LP-201 specifies
the ZAP universal wire used for the receipt encoding. LP-103 and LP-104 specify the explorer and
graph engine that index fills from the venue's trade rows.

\begin{thebibliography}{9}
\bibitem{lp032} Lux Industries, ``LP-032: DEX VM,'' Lux Proposals.
\bibitem{lp9999} Lux Industries, ``LP-9999: \texttt{0x9999} Receipt-Settlement Precompile,'' Lux Proposals.
\bibitem{lp305} Z.~Kelling, ``LP-311: One-Venue DEX Settlement,'' Lux Proposals, 2026.
\bibitem{lp020} Lux Industries, ``LP-020: Quasar Consensus Certificate,'' Lux Proposals.
\bibitem{lp073} Lux Industries, ``LP-073: Pulsar Post-Quantum Finality,'' Lux Proposals.
\bibitem{eip158} V.~Buterin, ``EIP-158: State clearing,'' Ethereum Improvement Proposals.
\end{thebibliography}

\end{document}
